\documentclass[11pt]{article}

\usepackage[margin=1in]{geometry}
\usepackage{amsmath,amssymb,amsfonts}
\usepackage{graphicx}
\usepackage{booktabs}
\usepackage{hyperref}
\usepackage{cleveref}
\usepackage{xcolor}

\title{\textbf{ACIES: Adaptive Perception Control \\
as Sequential Decision-Making over Perception Channels}}

\author{
  NICE-DEV226 \\
  \texttt{nicedev226@gmail.com}
}

\date{September 2026}

\begin{document}

\maketitle

\begin{abstract}
Deployed vision systems usually run their model at one fixed resolution, although many inputs
could be decided at a fraction of the cost. We treat the choice of how much perception to spend on
an input as a sequential decision problem. Each perception action (a resolution) is modelled as a
\emph{channel}: the quantised output of the model has a class-conditional distribution learned from
labelled examples. A controller keeps a belief about the decision and, by dynamic programming over
that belief, runs the next action only when its expected error reduction, priced by a single
parameter, exceeds its computational cost. We evaluate on real detections of YOLOv8n on COCO
val2017 with measured latency and an honest protocol (all settings chosen on a calibration split,
paired-bootstrap intervals on a held-out split). On the decision ``is there an object of class $c$
in the image?'' the controller uses 40\% less compute than fixed 640\,px inference for 0.2
accuracy points less on average over eight classes. Our first formulation, which described each
action by one independent ``clarity'' number, was \emph{worse than fixed-resolution inference on
every class}; we explain why. We also report two negative results: a simple two-stage confidence
cascade matches the cost/accuracy frontier, and no per-image resolution selector we tried improves
detection mAP over fixed-resolution YOLOv8n.
\end{abstract}

% ============================================================
\section{Introduction}
\label{sec:intro}

A clear scene and a cluttered one receive the same computation in a fixed-resolution pipeline.
Whether an input needs more perception is a question about \emph{that input}, and answering it
costs compute too. We ask when a controller can profit from deciding it online.

\textbf{Contributions.}
\begin{itemize}
  \item A formulation in which each perception action is a learned channel
  $P(\text{outcome}\mid Y,a)$ over a quantised model output, replacing a scalar ``clarity'', and a
  planner that trades expected error against measured cost (\Cref{sec:method}).
  \item A diagnosis, with measurements, of why the scalar-clarity formulation fails on real
  perception: repeated deterministic observations counted as independent, hard $0/1$ votes, and a
  safety override that fired on the uninformed prior (\Cref{sec:failure}).
  \item A reproducible real-model evaluation with a calibration/test protocol and confidence
  intervals, including negative results (\Cref{sec:experiments}).
\end{itemize}

% ============================================================
\section{Related Work}
\label{sec:related}

\textbf{Adaptive computation.} Early exiting \cite{teerapakayan2021early} and dynamic networks
\cite{han2021deep} adapt computation inside a model. We instead control \emph{which input
transformation} (resolution) a black-box model receives, so the approach applies to a model we
cannot modify.

\textbf{Cost-aware inference.} Cost-sensitive learning \cite{zadrozny2003cost} fixes costs at
training time; we choose actions at inference time from measured costs.

\textbf{Bandits and calibration.} Thompson sampling \cite{thompson1933likelihood,russo2018tutorial}
learns action quality online; we learn channels offline with Dirichlet posteriors and fit a
single evidence-discount parameter on held-out data, in the spirit of calibration
\cite{guo2017calibration}.

\textbf{Change-point detection.} Bayesian online change-point detection \cite{adams2007bayesian}
was part of an earlier version of the controller. With a constant hazard rate $h$ the posterior
mass of a change point at each step equals $h$ identically, so thresholding it cannot fire; we
found this by experiment and do not use change-point detection in the evaluated controller.

% ============================================================
\section{Method}
\label{sec:method}

\subsection{Setting}
Let $Y\in\{0,1\}$ be the decision variable with prior $\pi=P(Y{=}1)$. There are $n$ actions
$a_1,\dots,a_n$ ordered by increasing cost $C(a_i)$ (measured latency). Running $a$ on the input
yields a model output that is quantised into one of $M$ outcomes $o\in\{0,\dots,M{-}1\}$. Perception
is deterministic: running the same action again returns the same outcome.

\subsection{Channels}
Each action is described by two categorical distributions $P(o\mid Y{=}y,a)$, estimated from
labelled examples with a symmetric Dirichlet prior:
\begin{equation}
  \hat P(o\mid y,a)=\frac{n_{a,y,o}+\alpha}{\sum_{o'}n_{a,y,o'}+M\alpha}.
\end{equation}
The class prior $\pi$ is estimated from the same data. For $M=2$ and a symmetric table this is the
binary-symmetric channel of the original formulation.

\subsection{Belief with correlated actions}
With belief $b=P(Y{=}1\mid\text{history})$ and $\ell_a(o)=\log\frac{P(o\mid 1,a)}{P(o\mid 0,a)}$,
\begin{equation}
  \operatorname{logit} b' = \kappa\,\operatorname{logit} b + (1-\kappa)\operatorname{logit}\pi + \ell_a(o),
  \label{eq:update}
\end{equation}
where $\kappa\in[0,1]$ is the share of previous evidence that is kept. $\kappa{=}1$ is the exact
Bayes update for actions whose outputs are conditionally independent given $Y$; $\kappa{=}0$ lets
the latest outcome replace earlier ones. Errors of different resolutions on the same image are
correlated, so $\kappa{=}1$ counts the same difficulty twice. We fit $\kappa$ by held-out
log-loss on the calibration data.

\subsection{Planning}
Let $\lambda$ be the price of one error in cost units. Since re-running an action adds no
information and a cheaper action after a dearer one rarely does, plans move upward through the
ordered actions, and the state is $(b,j)$ where $j$ is the lowest action still allowed:
\begin{equation}
  V_t(b,j)=\min\Big\{\lambda\min(b,1{-}b),\ \min_{a_i,\,i\ge j}\Big[C(a_i)+\sum_{o}P(o\mid b,a_i)\,V_{t+1}(b'_{o},\,i{+}1)\Big]\Big\},
  \label{eq:bellman}
\end{equation}
with $V_H(b,j)=\lambda\min(b,1{-}b)$ for horizon $H$. We solve \Cref{eq:bellman} on a belief grid
with linear interpolation. The controller stops when stopping attains the minimum and decides
$\hat Y=\mathbf 1[b\ge\tfrac12]$. $\lambda$ is the single knob that trades accuracy for compute.
A greedy alternative (run the action maximising expected error reduction per unit cost while it
exceeds its price) is available; on a synthetic environment with a cheap uninformative first action
the plan attains 25\% lower expected objective (cost plus $\lambda\times$error) than the greedy rule.

\subsection{Exact optimum for the binary model}
For the binary-symmetric model with independent observations, the same recursion with $M=2$
and repeated actions allowed gives the exact optimal stopping policy; we use it as a lower bound
on cost against which any controller in that model can be measured (\texttt{acies.optimal}).

% ============================================================
\section{Why the scalar-clarity formulation fails}
\label{sec:failure}

The first version scored each action by $\Delta R/C$ with a single number $p_a=P(\text{correct}\mid a)$,
learned by Thompson sampling, and updated the belief as if every observation were an independent draw
of a symmetric channel. Three defects were measured.

\textbf{Emergency override on the prior.} The risk was scaled to $10\min(b,1{-}b)$, so the untouched
prior has risk $5.0$ above the emergency threshold $4.0$. The override, which forces the most
informative action irrespective of cost, therefore fired on every task and selected the most expensive
resolution: on 3000 simulated tasks the first action was the highest resolution 3000 times, and the
whole scoring machinery was bypassed. Restricting the override to states after at least one observation
halved both cost and error at confidence threshold 0.92 on the same simulator (cost 187 to 96, error 7.1\% to 3.3\%),
and cut cost by 39\% at equal error at threshold 0.99.

\textbf{Independent observations.} With deterministic perception, repeating an action returns the same
answer. The controller counted it as new evidence and paid again. On real YOLO detections this
appeared as the cheapest action being selected repeatedly.

\textbf{Hard votes.} A binary vote discards the detector's confidence, whereas the middle of the
confidence range is precisely where a second, costlier look pays off.

\textbf{Distance to the optimum.} On a simulated clarity table, the exact optimum of
\Cref{sec:method} costs 37 at 1.6\% error where the original controller cost about 240 (5--6$\times$ or
more); after the corrections above the gap is 3--4$\times$.

% ============================================================
\section{Experiments}
\label{sec:experiments}

\subsection{Setup}
\textbf{Data.} 3000 images sampled with a fixed seed from COCO val2017 (disjoint from the
train2017 images on which the Ultralytics weights are trained). \textbf{Model.} YOLOv8n, run at
$\{160,224,320,416,512,640,800,1024\}$\,px with confidence floor $0.001$ and NMS IoU $0.7$ (the
Ultralytics validation protocol); every detection and every latency is stored, so all policies are
evaluated on real outputs, none simulated. \textbf{Cost.} Latency on an 8-core CPU (4 threads),
per-resolution mean from a run on an idle machine ($22$--$183$\,ms). \textbf{Task.} For each of
eight classes, the binary decision ``is there an object of this class in the image?''; the outcome
is the top confidence of the class, quantised into nine bins. \textbf{Protocol.} Images are split
by index parity: even indices are calibration, odd indices are test (1500 images). Every setting is
chosen on calibration only, with one rule shared by all methods: the cheapest setting whose
cross-fitted calibration accuracy is within 0.3 points of fixed 640\,px. The chosen setting is then
frozen and measured once on the test split. Differences are given with paired-bootstrap 95\%
intervals over test images.

\textbf{Baselines.} Fixed resolution; a two-stage confidence cascade (run a cheap resolution, accept
if the confidence is above/below two thresholds, otherwise run a second resolution; resolutions and
thresholds selected on calibration by the same rule); and the original scalar-clarity controller.

\subsection{Results}
\Cref{tab:main} gives the test results. The original controller reached 89.2\%, 92.1\%, 92.4\%
and 93.9\% on \emph{person}, \emph{car}, \emph{chair} and \emph{bottle} at 67--78\,ms, below fixed
inference at lower cost on each. The channel controller (ACIES v2) saves 40\% of the compute of
fixed 640\,px on average for 0.2 accuracy points. Its loss is statistically significant on
\emph{person} ($-1.33$ points, 95\% CI $[-2.33,-0.40]$), borderline on \emph{cup}, and within the
noise on the other classes. The two-stage cascade attains a comparable or better frontier
(mean 95.6\% at 36\,ms).

\begin{table}[h]
\centering
\caption{Test split (1500 images per class). Accuracy and mean latency; $\Delta$ is the accuracy
difference of ACIES v2 from fixed 640\,px with a paired-bootstrap 95\% interval.}
\label{tab:main}
\begin{tabular}{lcccc}
\toprule
\textbf{Class} & \textbf{Fixed 640} & \textbf{ACIES v2} & $\Delta$ \textbf{[95\% CI]} & \textbf{Cascade} \\
\midrule
person & 94.7\% / 84\,ms & 93.4\% / 51\,ms & $-1.33$ [$-2.33$, $-0.40$] & 93.9\% / 45\,ms \\
car    & 96.7\% / 84\,ms & 96.5\% / 63\,ms & $-0.20$ [$-0.93$, $+0.47$] & 95.6\% / 44\,ms \\
chair  & 94.3\% / 84\,ms & 94.4\% / 46\,ms & $+0.14$ [$-0.73$, $+1.07$] & 93.7\% / 36\,ms \\
bottle & 95.6\% / 84\,ms & 95.9\% / 86\,ms & $+0.27$ [$-0.47$, $+1.00$] & 95.3\% / 46\,ms \\
cup    & 95.1\% / 84\,ms & 94.2\% / 32\,ms & $-0.95$ [$-1.93$, $0.00$]  & 94.8\% / 36\,ms \\
dog    & 97.6\% / 84\,ms & 98.0\% / 33\,ms & $+0.41$ [$-0.27$, $+1.07$] & 97.2\% / 31\,ms \\
tv     & 98.4\% / 84\,ms & 98.5\% / 53\,ms & $+0.15$ [$-0.40$, $+0.67$] & 98.3\% / 28\,ms \\
bench  & 96.7\% / 84\,ms & 96.2\% / 41\,ms & $-0.53$ [$-1.27$, $+0.20$] & 96.1\% / 24\,ms \\
\midrule
\textbf{Mean} & 96.1\% / 84\,ms & 95.9\% / 51\,ms & & 95.6\% / 36\,ms \\
\bottomrule
\end{tabular}
\end{table}

For \emph{bottle} the calibration rule selected the largest price in the grid and the controller
saved nothing (86\,ms): the rule can fail to find a cheaper setting.

\subsection{Negative result: detection quality}
We also asked whether a per-image resolution choice improves mean average precision (mAP50-95).
A selector spends one cheap pass, extracts label-free features of the detections (counts, confidence
mass, number of borderline detections, object sizes) and of the image (sharpness, contrast), predicts
per-resolution quality by ridge regression or boosted trees, and picks a resolution by predicted
quality minus a cost penalty. It does \emph{not} beat fixed-resolution YOLOv8n: reaching the mAP of
fixed 640\,px costs 101--124\,ms against 84\,ms. mAP of this model saturates near 512\,px, so the
first pass consumes the saving. In a 500-image pilot, an oracle that knows the ground truth would have saved about 60\%
at comparable quality, so the gap is one of prediction, not of headroom (the oracle
partly exploits label noise and is not attainable). In the same pilot the selector
seemed to beat the best fixed resolution by 1.1 mAP; at 1500 images the effect disappeared, which illustrates the size of the
sampling error at this scale.

\subsection{Limitations}
One detector (YOLOv8n) on one CPU, so cost ratios will differ on GPU or edge hardware; binary
presence decisions; 1500 test images per class (accuracy standard error about 0.6 points); COCO
only. The channel controller is not novel in its ingredients; its value here is a clean model of what
perception returns and an evaluation that separates it from a plain cascade. Not addressed:
multi-class decisions, temporal reuse across video frames, and online adaptation to distribution
shift.

% ============================================================
\section{Implementation and Reproducibility}
The controller, the exact binary optimum, the selector and the benchmark are open source
(Python, standard library only for the controller; a Go port and a C++ library implement the classic
simulator's primitives and are not part of the evaluated controller). The C++ library accessed
through \texttt{ctypes} was measured at 0.75--0.8$\times$ the speed of pure Python. All numbers in
this paper are produced by \texttt{benchmarks/final\_eval.py} from cached real detections; model
weights are not redistributed.

% ============================================================
\section{Conclusion}
Controlling perception per input is worth doing when the perception model's output is used as
evidence and the controller plans its escalation, and not worth it when it is described by one
independent number per action. On YOLOv8n it saves about 40\% of the compute for a decision
with a small accuracy cost, but a tuned cascade does as well, and it does not improve detection
mAP. We release the negative results with the evidence.

\bibliographystyle{plain}
\bibliography{references}

\end{document}
